% Chapter 10: Prepositions: Words of Position and Time

\chapter{Prepositions: Words of Position and Time}

\section{Pedagogical Purpose}
After learning about words that name (nouns), describe (adjectives), and act (verbs), students need to understand how these elements relate to each other in terms of position, time, and direction. Prepositions are crucial for building complex and meaningful sentences, moving beyond simple statements.

\begin{learningobjectives}
The learner can:
\begin{itemize}
    \item define what a preposition is.
    \item identify prepositions in sentences.
    \item categorize prepositions as showing place, time, or direction.
    \item use common prepositions correctly in sentences.
    \item understand that prepositions connect nouns/pronouns to other words.
\end{itemize}
\end{learningobjectives}

\section{Prerequisite Check}
\begin{quickcheck}
Underline the nouns and circle the verbs in each sentence.
\begin{enumerate}
    \item The \underline{dog} \textcircled{barks} loudly.
    \item \underline{Maya} \textcircled{is} happy.
    \item \underline{He} \textcircled{runs} fast.
\end{enumerate}
\end{quickcheck}

\section{Language Experience (Let's Begin)}
Ms. Sharma is giving instructions to Maya and Rohan for a treasure hunt.

\textbf{Ms. Sharma:} "Your first clue is hidden \textbf{under} the big tree. After you find it, look \textbf{behind} the swing. The next clue will be \textbf{on} the bench. You need to finish \textbf{by} lunchtime!"

\textbf{Rohan:} "So, we look \textbf{under} the tree, then \textbf{behind} the swing, and \textbf{on} the bench?"

\textbf{Maya:} "And we have to be done \textbf{by} lunchtime!"

\textbf{Ms. Sharma:} "Exactly! The words 'under', 'behind', 'on', and 'by' are very important. They tell us about the position of things or when something needs to happen. These words are called \textbf{Prepositions}."

\section{Concept Discovery (What is a Preposition?)}
A \textbf{preposition} is a word that shows the relationship between a noun or pronoun and other words in a sentence.

Prepositions often tell us about:
\begin{itemize}
    \item \textbf{Position/Place:} \textit{where} something is (e.g., \texttt{in}, \texttt{on}, \texttt{under})
    \item \textbf{Time:} \textit{when} something happens (e.g., \texttt{at}, \texttt{on}, \texttt{by})
    \item \textbf{Direction/Movement:} \textit{where} something is going (e.g., \texttt{to}, \texttt{into}, \texttt{from})
\end{itemize}

\begin{examplebox}
    The book is \textbf{on} the table.
    \begin{itemize}
        \item \texttt{on} shows the relationship between \texttt{book} and \texttt{table}.
    \end{itemize}
\end{examplebox}

\section{Concept Explanation (Kinds of Prepositions)}

\subsection*{A. Prepositions of Place}
\textbf{Prepositions of Place} tell us \textit{where} something is located.
\begin{itemize}
    \item The cat is \textbf{under} the bed.
    \item The picture is \textbf{on} the wall.
    \item Rohan is sitting \textbf{between} Maya and David.
\end{itemize}

\subsection*{B. Prepositions of Time}
\textbf{Prepositions of Time} tell us \textit{when} something happens.
\begin{itemize}
    \item We have school \textbf{on} Monday.
    \item The movie starts \textbf{at} 7 o'clock.
    \item I will finish my homework \textbf{by} evening.
\end{itemize}

\subsection*{C. Prepositions of Direction}
\textbf{Prepositions of Direction} tell us \textit{where} something is moving.
\begin{itemize}
    \item The car went \textbf{through} the tunnel.
    \item Rohan is walking \textbf{to} the school.
    \item The cat jumped \textbf{over} the wall.
\end{itemize}

\section{Guided Practice (Let's Practise Together)}
\begin{activitybox}{Activity A: Find the Preposition and its Kind}
\textbf{Ms. Sharma:} "Let's find the prepositions in these sentences and name their type. I'll do the first one."
\begin{enumerate}
    \item The book is \textbf{on} the table.
    \begin{itemize}
        \item \textit{Preposition:} \texttt{on}
        \item \textit{Kind:} Preposition of Place
    \end{itemize}
    \item The school starts \textbf{at} 9 AM.
    \item The dog ran \textbf{across} the road.
    \item My birthday is \textbf{in} January.
\end{enumerate}
\end{activitybox}

\section{Summary}
\begin{summarybox}
In this chapter, we learned that \textbf{prepositions} are connecting words that show relationships.
\begin{itemize}
    \item They often tell us about position, time, or direction.
    \item \textbf{Preposition of Place}: Tells \textit{where} something is. (e.g., The cat is \textbf{on} the roof.)
    \item \textbf{Preposition of Time}: Tells \textit{when} something happens. (e.g., The class starts \textbf{at} 9 AM.)
    \item \textbf{Preposition of Direction}: Tells \textit{where} something is going. (e.g., He walked \textbf{to} the store.)
\end{itemize}
Prepositions help us build clear and detailed sentences by showing how things are related!
\end{summarybox}